This paper draws on four literatures: prediction-market prices as probabilities, models of
probability-price dynamics, option pricing and the volatility smile, and reduced-form credit
models.

\paragraph{Prediction-market prices as probabilities.} \citet{wolferszitzewitz2004} survey prediction
markets and their forecasting record. \citet{manski2006} and \citet{wolferszitzewitz2006} ask when
an equilibrium price equals a mean belief. \citet{berg2008} document long-run accuracy in the Iowa
Electronic Markets, and \citet{pageclemen2013} find that prediction markets are well calibrated
except for contracts with long horizons. In betting markets the dominant departure from
calibration is the favourite--longshot bias: longshots are overpriced relative to their win
frequency \citep{griffith1949,thalerziemba1988}. \citet{snowbergwolfers2010} attribute it to
misperception of small probabilities rather than risk love. \citet{ottaviani2008} survey the
explanations. On Kalshi, \citet{burgi2025} use transaction data on more than 300,000 contracts.
They document a favourite--longshot bias that is stronger for takers than for makers, and model it
as modest belief disagreement with an overestimation of small probabilities.
\citet{rothschildsethi2016} study trader types and microstructure in a political prediction market.

These papers study the \emph{level} of prices relative to outcomes. We study how prices
\emph{move}. We also show that the favourite--longshot bias reappears, in different guises, in the
cross-section of strikes and in the term structure of deadlines.

\paragraph{Dynamics of probability prices.} \citet{stern1994} models the score difference of a game
as a Brownian motion. The win probability is then a normal distribution function of the lead,
scaled by the square root of time remaining: our M0 with a uniform clock.
\citet{polsonstern2015} invert betting lines to obtain the implied volatility of a game, and
\citet{fengpolsonxu2016} extend the approach to soccer with a Skellam process.

\citet{archakipeirotis2009} derive the volatility of a binary contract's price from latent ability
diffusions. They show that volatility depends on the price and on time to expiration, and that the
structural model beats GARCH in forecasting InTrade volatility. \citet{taleb2018} uses the same
construction to argue that a well-calibrated election forecast must stay close to one half when
uncertainty is high.

Two recent papers are the closest to ours.
\begin{itemize}
  \item \citet{dalen2025} proposes a ``Black--Scholes for prediction markets''. It models log-odds as
  a jump-diffusion whose drift is fixed by the martingale condition. It describes an EM
  calibration that separates diffusion from jumps, and a layer of variance and correlation
  derivatives. Its validation is largely on synthetic paths.
  \item \citet{xi2026} estimate a structural volatility model on a large panel of Kalshi contracts.
  It combines a Wright--Fisher deadline-resolution component with a Glosten--Milgrom order-flow
  component \citep{glostenmilgrom1985}. They find that it forecasts volatility better than ARCH and
  GARCH benchmarks, that volatility peaks near even odds and rises into resolution, and that sports
  contracts jump while economics contracts move smoothly.
\end{itemize}
\citet{feilnendel2026} derive optimal market making when the price is a transformed latent-belief
diffusion; inventory-based market making goes back to \citet{avellanedastoikov2008}. Cost-function
market makers \citep{hanson2003} are an alternative price-formation mechanism, which Kalshi does
not use.

Relative to this work, we make five contributions.
\begin{itemize}
  \item We state the exact transition law of the Gaussian model and test it as a distributional
  restriction: the standardised innovation $\ustd$ must be $\N(0,1)$ at every horizon, given the
  clock.
  \item We estimate information clocks, backbones, stochastic volatility and jumps by contract
  category on out-of-sample data. At high frequency, we use simulated controls that reproduce the
  tick grid.
  \item We extend the framework to strips of contracts on a numeric outcome. There, implied
  distributions, smiles, smile dynamics and a variance-swap analogue are observable.
  \item We treat deadline contracts as a reduced-form hazard model, with implied hazard term
  structures, theta and event timing.
  \item We show that a common tail-and-time premium links these settings.
\end{itemize}
Our estimated backbone is close to the Wright--Fisher exponent of one half. This independently
supports the specification of \citet{xi2026}.

\paragraph{Option pricing, smiles and variance premia.} The model is a cash-or-nothing digital
option under Bachelier dynamics
\citep{bachelier1900,blackscholes1973,merton1973}. The extensions follow familiar steps:
\begin{itemize}
  \item jumps \citep{merton1976};
  \item stochastic volatility \citep{heston1993};
  \item local and stochastic smiles \citep{dupire1994,hagan2002}.
\end{itemize}
SABR's backbone exponent and correlation have direct counterparts in M1.
\citet{breedenlitzenberger1978} show that a strip of option prices identifies the risk-neutral
density; a strip of digitals does so directly. The distinction between sticky-strike and
sticky-delta smile dynamics is due to \citet{derman1999}; \citet{gatheral2006} reviews smile
modelling. Comparing realized with implied variance follows the variance-risk-premium literature
\citep{carrwu2009,bollerslev2009}.

\paragraph{Reduced-form credit.} A contract that pays if an event occurs by a deadline is the
complement of a defaultable zero-coupon claim with zero recovery. Intensity models
\citep{jarrowturnbull1995,lando1998,duffiesingleton1999} supply the implied hazard, forward
hazards and the compensated-jump martingale. The square-root intensity of \citet{cox1985cir} is
the natural Cox specification.

\paragraph{Econometrics of volatility and jumps.} At high frequency we use:
\begin{itemize}
  \item volatility signature plots \citep{andersen2000signature} and the noise framework of
  \citet{zhang2005}, to separate microstructure noise from volatility;
  \item bipower variation \citep{barndorffnielsen2004} and threshold tests
  \citep{leemykland2008}, to separate jumps from diffusion;
  \item a bid--ask bounce measure \citep{roll1984} and price discreteness \citep{harris1991}, for
  the tick grid.
\end{itemize}
The EM decomposition of hourly moves into diffusion and jumps follows \citet{dempster1977}.
